\documentclass[11pt]{article}
\usepackage[margin=1in]{geometry}
\usepackage{amsmath, amssymb}
\usepackage{enumitem}
\usepackage{titlesec}
\usepackage{hyperref}
\usepackage{booktabs}
\usepackage{xcolor}
\usepackage{parskip}

\hypersetup{colorlinks=true, linkcolor=blue, urlcolor=blue}

\title{\textbf{EE627 --- Class Summary}\\[0.4em]
       \large February 27, 2026}
\author{Prof.\ Rensheng Wang}
\date{}

\begin{document}
\maketitle
\tableofcontents
\newpage

%------------------------------------------------------------
\section{Final Project Structure and Course Philosophy}
%------------------------------------------------------------

\subsection{How the Final Project Works}
Starting this week, the course pivots to an industry recommender dataset.
The project is \textbf{guided} --- not ``here is data, figure it out.''
\begin{itemize}
    \item Each week introduces one new technique or approach.
    \item Sample code and a direction are provided.
    \item Students apply that week's technique to the same ongoing project.
    \item By the end of the semester, multiple approaches have been tried
      on the same dataset, allowing direct comparison.
\end{itemize}
The goal is to replicate the workflow of a real data scientist iterating
toward a better solution --- not to complete a one-shot assignment.

\subsection{Why This Class Exists}
The professor explained the origin of the course directly:
\begin{itemize}
    \item ECE had no recommender-focused course despite recommenders being
      one of the highest-cash-flow areas in tech.
    \item The syllabus is assembled from \textbf{industry interview
      requirements}, not from textbooks. Every topic covered reflects
      what top companies actually test.
    \item The interview process at major tech companies has become highly
      standardized --- there are known preparation paths for specific
      companies (behavioral, system design, ML technical). This course
      targets the ML technical component.
\end{itemize}

\subsection{The Right Way to Engage}
\begin{itemize}
    \item \textbf{Ask questions aggressively.} ``What happens if my data
      is NOT like what you described?'' is the right question. The
      standard case is never what you encounter in a real job.
    \item \textbf{Challenge yourself.} After every explanation: does this
      make sense? Can I generalize it beyond the example shown?
    \item Focus on \textbf{higher-level reasoning} --- why to choose model
      A over B, what assumptions to check, how to handle dirty data ---
      not on syntax and commands.
\end{itemize}


\newpage
%------------------------------------------------------------
\section{Two-Tower Model: Training and Tuning}
%------------------------------------------------------------

This lecture provided the deepest explanation of how the two-tower
recommendation model actually works end-to-end, including training.

\subsection{The Core Idea: Everything Becomes a Number}
Computers cannot process names, IDs, or descriptions. The fundamental
operation of the two-tower system is:

\[
\textbf{user features} \;\xrightarrow{\text{encoder}}\; \mathbf{u} \in
\mathbb{R}^d, \qquad
\textbf{item features} \;\xrightarrow{\text{encoder}}\; \mathbf{v} \in
\mathbb{R}^d
\]
\[
\text{relevance score}(u, i) = \mathbf{u} \cdot \mathbf{v}
\]

\begin{itemize}
    \item A user named ``Jason'' with membership ID ``12700-AD-Y'' gets
      mapped to a float vector, e.g., $[0.3,\; -0.4,\; -0.1,\; \ldots]$.
    \item A video about deep learning gets its own float vector.
    \item Dot product between the two vectors $=$ a single scalar
      representing predicted affinity.
    \item \textbf{High dot product} $\Rightarrow$ user likely to engage
      with this item.
    \item \textbf{Low dot product} $\Rightarrow$ user unlikely to engage.
\end{itemize}

\subsection{Training: From Random to Meaningful}
\begin{enumerate}
    \item \textbf{Initialization:} assign random float vectors to every
      user and every item. At this point, the dot products are meaningless.
    \item \textbf{Observe interactions:} the system collects real user
      behavior --- clicks, purchases, watch time, shares.
    \item \textbf{Tune embeddings:} adjust the vectors so that user--item
      pairs with high interaction get high dot products, and pairs with
      no interaction get low dot products.
    \item \textbf{Iterate:} after every batch of new interactions,
      re-tune. The embeddings gradually become meaningful representations.
\end{enumerate}

\textit{Intuition:} if a user watches many deep learning videos, their
embedding vector should drift toward the region of the latent space where
deep learning video vectors cluster. ``Closeness in the embedding space
$\Leftrightarrow$ high likelihood of engagement.''

\subsection{Offline vs.\ Online: Why the Split Is Necessary}
Training embeddings is \textbf{computationally expensive} --- it cannot
happen in real time.

\begin{center}
\begin{tabular}{@{}lll@{}}
\toprule
\textbf{Pipeline} & \textbf{Timing} & \textbf{Work done} \\
\midrule
Offline & Nightly (or periodic batch) & Retrain/update embeddings using all interactions \\
Online  & Milliseconds (user request) & Fetch pre-trained embeddings; compute dot products; rank \\
\bottomrule
\end{tabular}
\end{center}

\begin{itemize}
    \item When a user logs in to Amazon at 9 a.m., the system \emph{fetches}
      their pre-computed embedding and computes dot products against item
      embeddings instantly.
    \item The embedding models are \emph{not} retrained at login time.
    \item After midnight, the system processes the day's interactions
      offline to update the embeddings for the next day.
\end{itemize}


\newpage
%------------------------------------------------------------
\section{Retrieval (Stage 1): Vector Nearest-Neighbor Search}
%------------------------------------------------------------

\subsection{The Geometric Picture}
Think of the embedding space as a high-dimensional map. Each item is a
point on this map. Each user query is also a point. The Stage 1 retrieval
problem is: \textit{find the points closest to the user query point.}

\begin{itemize}
    \item Items nearby in the embedding space have high dot products with
      the user vector.
    \item Items far away have low dot products.
    \item A ``slice'' of the nearest neighborhood gives the candidate pool.
\end{itemize}

\textit{Analogy from lecture:} imagine a pie chart. The user vector is the
center. Items form a ring around it. The retrieval algorithm takes the
nearest arc --- the items within the smallest angular distance from the
user vector.

\subsection{Why ANN (Approximate Nearest Neighbor) Is Needed}
A billion-item catalog means a billion dot products per user request.
Computing all of them exactly is too slow. \textbf{Approximate Nearest
Neighbor (ANN)} algorithms (e.g., HNSW, FAISS) trade a small accuracy
loss for orders-of-magnitude speed improvement, reducing lookup time
from seconds to milliseconds.

\subsection{Output}
Stage 1 outputs $\sim$100--1000 candidate items per user.


\newpage
%------------------------------------------------------------
\section{Feature Engineering for Stage 3 Fine Ranking}
%------------------------------------------------------------

Stage 3 is the \textbf{sophisticated feature engineering stage}. Because
the candidate pool is now small (hundreds of items), a rich ML model
can be run on each (user, item) pair. The quality of the features
determines the quality of the ranking.

\subsection{User-Side Features}
\begin{itemize}
    \item \textbf{User ID:} the raw identifier, embedded into a vector.
    \item \textbf{Login time:} hour of day (6 a.m.\ vs.\ 10 p.m.\ reflect
      different user states and intent). Encoded as one-hot or cyclic
      numeric features.
    \item \textbf{Region / location:} inferred from IP address.
      Enables local recommendations (nearby news, local events, relevant
      language content).
    \item \textbf{Language / device:} proxy for nationality, platform
      preference, and content type.
    \item \textbf{Historical behavior (statistics):} the richest signal.
      In the past week:
      \begin{itemize}
          \item Number of purchases / views by category
          \item Distribution across categories (electronics: 3, fashion: 4,
            sports: 2)
          \item Average session length, frequency
          \item Which specific items were browsed, purchased, reviewed
      \end{itemize}
\end{itemize}

\subsection{Item-Side Features}
\begin{itemize}
    \item Item ID (embedded vector)
    \item Click-through rate (CTR), watch time, share count
    \item Category, tags, description keywords
    \item Recency (how recently was this published/listed?)
    \item Popularity rank within category
\end{itemize}

\subsection{Model Choices for Stage 3}
All of these features feed into a standard ML model:
\begin{itemize}
    \item XGBoost / CatBoost / GBDT (tabular features)
    \item Deep neural network (mix of sparse and dense features)
    \item LLM-based ranker (emerging)
\end{itemize}
The model predicts an engagement probability score for each
(user, item) pair. Items are sorted by score.


\newpage
%------------------------------------------------------------
\section{Re-ranking (Stage 4): Beyond Pure Relevance}
%------------------------------------------------------------

Stage 4 is where business, psychology, and ethics enter the pipeline.
The ranked list from Stage 3 is technically correct but not yet ready
to show to the user.

\subsection{The Dead Loop Problem (Filter Bubble)}
If the system only recommends what the user has already shown interest in:
\begin{enumerate}
    \item User watches gaming videos.
    \item System recommends more gaming videos.
    \item User watches more gaming videos.
    \item System recommends even more gaming videos.
    \item User never discovers anything else.
\end{enumerate}
This is a \textbf{filter bubble} --- the user's world shrinks to their
existing interests. It also limits revenue, because a user who only
consumes one type of content has a limited engagement ceiling.

\subsection{Diversity Injection}
Intentionally include a small fraction ($\sim$10\%) of items outside the
user's primary interests:
\begin{itemize}
    \item If the user primarily watches gaming: add 1--2 sports videos.
    \item If the user searches for swimwear: add 1--2 general outdoor
      sports items.
    \item Purpose: \textbf{spark new interests}, expand the user's
      engagement surface, and increase long-term revenue.
\end{itemize}
\textit{This is not a mistake and not random --- it is intentional
psychological engineering.}

\subsection{Business Rules (Sponsored Content)}
\begin{itemize}
    \item Advertisers pay for placement. Some items are elevated in the
      final ranking \emph{despite lower relevance scores} because a
      business contract requires it.
    \item Example: Amazon search results include sponsored products that
      are ranked by payment, not by relevance.
    \item The system must blend: (high relevance score) + (sponsored
      placement contract) into the final list.
\end{itemize}

\subsection{Contextual Adjustment}
The final ranking accounts for situational context:
\begin{itemize}
    \item \textbf{Timestamp / calendar:} on Thanksgiving, inject holiday
      gift recommendations regardless of the user's search intent.
    \item \textbf{Trending events:} a major news story justifies appearing
      in feeds even for users who did not search for it.
    \item \textbf{Freshness:} prefer recently published content over older
      content with equal relevance scores.
\end{itemize}

\subsection{Evaluation: How to Know If Re-ranking Is Working}
Change one parameter (e.g., diversity ratio from 4\% to 5\%) and measure:
\begin{itemize}
    \item Average session duration (did users spend more time?)
    \item Click-through rate on diversified items (did new interests spark?)
    \item Revenue change
\end{itemize}
At billion-user scale, a 1\% improvement in session duration translates
to hundreds of millions in additional revenue. This is why companies
invest heavily in re-ranking tuning and why engineers who do it well
command exceptional compensation.


\newpage
%------------------------------------------------------------
\section{Interview Preparation: Drawing the Recommender Pipeline}
%------------------------------------------------------------

The professor made explicit what industry interviews test for recommenders:

\textit{``They will not ask you to write code for a recommender. They will
ask you to \textbf{draw the diagram}.''} Specifically:

\begin{enumerate}
    \item Draw the four-stage pipeline (Retrieval $\to$ Rough Ranking
      $\to$ Fine Ranking $\to$ Re-ranking).
    \item Draw the two-tower architecture: user tower, item tower,
      dot product, interaction-based training signal.
    \item Draw the offline/online split: offline nightly training,
      online real-time serving.
    \item For each stage: state the \textbf{candidate pool size}, the
      \textbf{key technique}, and the \textbf{output size}.
    \item Explain \textbf{why} each stage exists --- not just what it does.
\end{enumerate}

\begin{center}
\begin{tabular}{@{}llll@{}}
\toprule
\textbf{Stage} & \textbf{Input} & \textbf{Technique} & \textbf{Output} \\
\midrule
1.\ Retrieval       & Billions & Two-tower ANN search & $\sim$1000 \\
2.\ Rough Ranking   & $\sim$1000 & Rules, shallow ML, Bloom filter & $\sim$300 \\
3.\ Fine Ranking    & $\sim$300 & Feature engineering + XGBoost/DNN & $\sim$50 \\
4.\ Re-ranking      & $\sim$50 & Business rules, diversity, context & Top 10 \\
\bottomrule
\end{tabular}
\end{center}


\newpage
%------------------------------------------------------------
\section{Final Project: Yahoo Music Dataset (Detailed)}
%------------------------------------------------------------

The course final project uses a real industry dataset from \textbf{Yahoo
Music}. The core task is building a recommendation system that predicts
user--item interactions (song ratings).

\subsection{Hierarchical Data Structure}
Music data is stored in a four-level hierarchy:

\[
\text{Track} \;\longrightarrow\; \text{Album} \;\longrightarrow\;
\text{Artist} \;\longrightarrow\; \text{Genre}_1, \ldots, \text{Genre}_K
\]

\begin{itemize}
    \item A song track belongs to one album.
    \item An album belongs to one artist.
    \item An artist is associated with one or more genres.
    \item Users can rate \emph{any} item in the hierarchy (track, album,
      artist, or genre) with a score between 0 and 100.
\end{itemize}

\textit{Example:} ``Poker Face'' $\to$ ``The Fame'' $\to$ ``Lady Gaga''
$\to$ Pop, Dance.

\subsection{Data Statistics}

\begin{center}
\begin{tabular}{@{}lll@{}}
\toprule
\textbf{Split} & \textbf{Quantity} & \textbf{Notes} \\
\midrule
Training users  & 249K & --- \\
Items           & 296K & tracks, albums, artists, genres \\
Rating scores   & 62M  & explicit ratings 0--100 \\
Test users      & 100K & subset of training users \\
Test tracks/user & 6   & must classify each as rated or not \\
\bottomrule
\end{tabular}
\end{center}

Training data item breakdown by level: tracks 43.6\%, albums 31.2\%,
artists 19.1\%, genres 6.0\%. Test data is \emph{100\% track-level}.

\subsection{The Actual Task (Binary Classification)}
For each test user, 6 candidate track IDs are provided. The task is
\textbf{not} to predict the rating score --- it is to determine which
3 of the 6 tracks the user actually rated highly (label ``1'') and which
3 they did not (label ``0'').

\begin{itemize}
    \item Training data: (UserID, ItemID, Rating) tuples across all item
      levels.
    \item Test data: (UserID, TrackID, ?) --- binary label unknown.
    \item Goal: produce correct 3-``1'' / 3-``0'' assignments.
\end{itemize}

\subsection{Hierarchy Search: Linking Tracks to Ratings}
A test track often has \emph{no direct rating} in the training data.
The solution is \textbf{hierarchy search}: look up the track's parent
album, parent artist, and associated genres and retrieve whatever
ratings exist at those levels.

\textit{Example for User 249008, Track 1 (ID 4967):}
\begin{itemize}
    \item Album 205719 was rated 90; Artist 197877 was rated 80.
    \item Genre IDs 172023, 201738 present but no direct genre rating.
\end{itemize}
The hierarchy ratings become the feature set for the prediction model.

\subsection{Feature Generation from Rating Scores}

\textbf{Approach 1 --- Simple rule-based:}
\begin{itemize}
    \item Use raw hierarchy ratings directly as a heuristic.
    \item Achieves $\sim$87\% accuracy (13\% error rate).
\end{itemize}

\textbf{Approach 2 --- Statistical features:}
\begin{itemize}
    \item Album and artist each have one rating score per user.
    \item Genres may have \emph{multiple} rating scores (user rated
      several tracks in the genre); extract:
      \begin{itemize}
          \item Maximum / minimum score
          \item Average score
          \item Variance
          \item Score count (length)
      \end{itemize}
    \item Feed these features into an ML classifier.
    \item Achieves $\sim$90\% accuracy (10\% error rate).
\end{itemize}

\subsection{Handling Sparse Users: Similar-User Borrowing}
Some users have very few rating records. For these users, hierarchy
search and expansion return nothing useful. Strategy:

\begin{enumerate}
    \item \textbf{Hierarchy expansion:} from a known item, expand
      upward through the hierarchy tree to find \emph{all} related
      items and borrow their ratings as proxy features.
    \item \textbf{Similar-user matching:} find training users with
      similar rating behavior; borrow their hierarchy features to
      supplement the sparse test user's feature vector.
\end{enumerate}

\subsection{Actual File Formats}

\textbf{Test data} --- \texttt{testTrack\_hierarchy.txt} (pipe-delimited):
\begin{center}
\texttt{userID | trackID | albumID | artistID | genreID$_1$ | $\cdots$ | genreID$_K$}
\end{center}
Fields that do not exist are filled with \texttt{None}.
Example rows:
\begin{verbatim}
199810|208019|209288|None
199810|74139|277282|271146|113360|173467|173655|192976|...
\end{verbatim}

\textbf{Training data} --- \texttt{trainIdx2\_matrix.txt} (pipe-delimited):
\begin{center}
\texttt{userID | itemID | score}
\end{center}
\texttt{itemID} may be a track, album, artist, or genre ID; \texttt{score}
is 0--100.
Example rows:
\begin{verbatim}
199808|248969|90
199808|2663|90
199808|28341|90
\end{verbatim}

The workflow: for each test user, retrieve their test track hierarchy from
\texttt{testTrack\_hierarchy.txt}, then look up each hierarchy node in
\texttt{trainIdx2\_matrix.txt} to collect the user's historical rating
scores as features.

Sample starter code is provided in \texttt{read\_rating\_V3.py}.

\subsection{The Approach: Matrix Factorization}
The first technique applied to the project will be
\textbf{matrix factorization} --- the foundational collaborative filtering
method. The goal: decompose the sparse rating matrix into compact user and
item embedding vectors such that the dot product of a user vector and an
item vector predicts the user's rating for that item.

This is precisely the two-tower concept applied to an explicit rating
matrix. The details and implementation will be covered in subsequent
lectures.

\subsection{Course Trajectory}
Each week, a new technique is introduced and applied to the same Yahoo
Music dataset:
\begin{enumerate}
    \item Matrix factorization (collaborative filtering baseline)
    \item Additional retrieval and ranking approaches
    \item Feature engineering to improve Stage 3 scoring
    \item Evaluation metrics (comparing models)
\end{enumerate}
The final grade reflects the student's ability to improve the recommender
across multiple iterations --- mirroring the real job of a data scientist.

\end{document}
